\section{Part C: Alignment and Limitations}
SmartTE agrees with the literature in one limited sense: path diversity can make it possible to avoid a local bottleneck when a controller has usable network-state information \cite{hedera,microte}. It does not reproduce the scale, workload knowledge, rate control, or operational guarantees of B4, SWAN, or large data-center schedulers \cite{b4,swan}. In this emulation, the static E2 references establish that moving the primary flow away from the Path-A competitor can preserve throughput.

The observed policy trade-off also matches the motivation for route-stability controls. B changes paths frequently; the frozen C setting makes fewer successful changes in E1--E3, but it may wait longer before acting. This is a behavior of the measured implementation, not proof that a particular margin/dwell pair is optimal. SmartTE does not provide the formal update-consistency guarantees studied in prior work \cite{reitblatt}; it records controller events and differentiates successful transitions from rejected decisions.

The interpretation is constrained by a single-host Mininet/OVS emulation under WSL2, a four-switch topology, two paths with unequal configured delays, one TCP pair, and one optional UDP competitor. Emulated timing, Linux queues, host load, and the recorded HTB warning can affect results. Queue limits were verified, but transient occupancy was not continuously measured. E4 is a diagnostic one-link failure: B had already moved to B before failure in all five trials, while C remained on A and changed after failure in all five. Each had one successful change per trial. Raw decision attempts averaged 17 for B and 11.8 for C; rejected attempts averaged 16 and 10.8, respectively. Static stayed on the failed path in all five trials. E4 therefore cannot rank failover latency fairly. The campaign metadata records commit 58acacd120a1b8b85b81d8a46c1f273479d18942 with \texttt{git\_dirty=true}, so the exact uncommitted source state is not fully established. The results characterize this code, topology, and workload rather than a production traffic-engineering system.
